\documentclass[conference]{IEEEtran}
\usepackage{cite}
\usepackage{url}
\usepackage{amsmath,amssymb}
\usepackage{graphicx}
\usepackage{booktabs}
\usepackage{pgfplots}
\usepgfplotslibrary{groupplots}
\pgfplotsset{compat=1.18}

% Generated by paper_ieee/scripts/make_artifacts.py; do not edit.
\newcommand{\AQHeldOutQuestions}{114}  % data/results/aq/sample.csv
\newcommand{\AQSampleShare}{20\%}  % scripts/aq_run.py at paper-v1.2
\newcommand{\AQScannedQuestions}{25}  % data/results/aq/sample.csv
\newcommand{\AQTuningQuestions}{179}  % data/results/aq/sample.csv
\newcommand{\AccountingHeldOut}{31}  % pdfs/split.csv
\newcommand{\AccountingMain}{17}  % data/results/p4/P4-FINAL-20261106-1_primary_documents.csv
\newcommand{\AccountingRemovedOwnHindiTrack}{2}  % pdfs/split.csv
\newcommand{\AccountingRemovedQASPER}{3}  % pdfs/split.csv
\newcommand{\AccountingRemovedUniDocBench}{9}  % pdfs/split.csv
\newcommand{\BMKone}{1.5}  % runs/P3-LOCAL-20261105-4.json
\newcommand{\BMb}{0.75}  % runs/P3-LOCAL-20261105-4.json
\newcommand{\BudgetHigh}{2{,}000}  % runs/P3-LOCAL-20261105-4.json
\newcommand{\BudgetLow}{750}  % runs/P3-LOCAL-20261105-4.json
\newcommand{\BudgetMain}{1{,}250}  % runs/P3-LOCAL-20261105-4.json
\newcommand{\CILevel}{95}  % runs/P4-FINAL-20261106-1.json
\newcommand{\ChunkMax}{1{,}024}  % src/chunk/base.py and src/chunk/chunkers.py and src/eval/stats.py and scripts/p4_analyse.py and scripts/aq_run.py and scripts/golden_check.py and scripts/p3_fixedk.py at paper-v1.2
\newcommand{\ChunkMin}{100}  % src/chunk/base.py and src/chunk/chunkers.py and src/eval/stats.py and scripts/p4_analyse.py and scripts/aq_run.py and scripts/golden_check.py and scripts/p3_fixedk.py at paper-v1.2
\newcommand{\CorrectionCount}{5}  % spec/spec_v1.2.md at paper-v1.2, rows of kind correction in tables/deviations.tex
\newcommand{\CoverageBySourceHeldOutCCUADDocling}{0\%}  % data/results/p3/gold_drop_final.csv
\newcommand{\CoverageBySourceHeldOutCCUADPdfplumber}{0\%}  % data/results/p3/gold_drop_final.csv
\newcommand{\CoverageBySourceHeldOutCOHRBenchDocling}{22\%}  % data/results/p3/gold_drop_final.csv
\newcommand{\CoverageBySourceHeldOutCOHRBenchPdfplumber}{37\%}  % data/results/p3/gold_drop_final.csv
\newcommand{\CoverageBySourceHeldOutCUADDocling}{0\%}  % data/results/p3/gold_drop_final.csv
\newcommand{\CoverageBySourceHeldOutCUADPdfplumber}{0\%}  % data/results/p3/gold_drop_final.csv
\newcommand{\CoverageBySourceHeldOutIndianDocling}{2\%}  % data/results/p3/gold_drop_final.csv
\newcommand{\CoverageBySourceHeldOutIndianPdfplumber}{14\%}  % data/results/p3/gold_drop_final.csv
\newcommand{\CoverageBySourceHeldOutOHRBenchDocling}{30\%}  % data/results/p3/gold_drop_final.csv
\newcommand{\CoverageBySourceHeldOutOHRBenchPdfplumber}{34\%}  % data/results/p3/gold_drop_final.csv
\newcommand{\CoverageBySourceHeldOutOodFinanceBenchDocling}{86\%}  % data/results/p3/gold_drop_final.csv
\newcommand{\CoverageBySourceHeldOutOodFinanceBenchPdfplumber}{86\%}  % data/results/p3/gold_drop_final.csv
\newcommand{\CoverageBySourceHeldOutOodOHRBenchDocling}{3\%}  % data/results/p3/gold_drop_final.csv
\newcommand{\CoverageBySourceHeldOutOodOHRBenchPdfplumber}{19\%}  % data/results/p3/gold_drop_final.csv
\newcommand{\CoverageHeldOutDoclingLost}{121}  % data/results/p3/gold_drop_final.csv
\newcommand{\CoverageHeldOutDoclingShare}{16.3\%}  % data/results/p3/gold_drop_final.csv
\newcommand{\CoverageHeldOutDoclingTotal}{742}  % data/results/p3/gold_drop_final.csv
\newcommand{\CoverageHeldOutPdfplumberLost}{205}  % data/results/p3/gold_drop_final.csv
\newcommand{\CoverageHeldOutPdfplumberShare}{27.6\%}  % data/results/p3/gold_drop_final.csv
\newcommand{\CoverageHeldOutPdfplumberTotal}{742}  % data/results/p3/gold_drop_final.csv
\newcommand{\CoverageTuningDoclingLost}{84}  % data/results/p3/gold_drop_tuning.csv
\newcommand{\CoverageTuningDoclingShare}{8.7\%}  % data/results/p3/gold_drop_tuning.csv
\newcommand{\CoverageTuningDoclingTotal}{970}  % data/results/p3/gold_drop_tuning.csv
\newcommand{\CoverageTuningPdfplumberLost}{92}  % data/results/p3/gold_drop_tuning.csv
\newcommand{\CoverageTuningPdfplumberShare}{9.5\%}  % data/results/p3/gold_drop_tuning.csv
\newcommand{\CoverageTuningPdfplumberTotal}{970}  % data/results/p3/gold_drop_tuning.csv
\newcommand{\DateHeldOutScored}{6 November 2026}  % runs/P3-FINAL-20261106-1.json created_at
\newcommand{\DateScorerFrozen}{5 November 2026}  % git 448173c, tag scorer_v1
\newcommand{\DateSpecVone}{26 October 2026}  % git 607d29c, first adds spec/spec_v1.md
\newcommand{\DateSpecVoneOne}{1 November 2026}  % git de4aff3, first adds spec/spec_v1.1.md
\newcommand{\DateSpecVoneTwo}{3 November 2026}  % git b1bd7d1, first adds spec/spec_v1.2.md
\newcommand{\DeviationCount}{21}  % spec/spec_v1.2.md at paper-v1.2, rows of kind deviation in tables/deviations.tex
\newcommand{\DoclingVersion}{2.143.0}  % runs/P1-RQ5-KN01-20261109-4.json
\newcommand{\EmbedDim}{1{,}024}  % runs/P3-LOCAL-20261105-4.json
\newcommand{\FixedKBudget}{4{,}000}  % runs/P3-FIXEDK-20261112-1.json
\newcommand{\GoldenIndianDFourRtiActInvalid}{8}  % data/results/golden_check/labels.csv
\newcommand{\GoldenIndianDFourRtiActLabelled}{20}  % data/results/golden_check/labels.csv
\newcommand{\GoldenIndianDFourRtiActRate}{0.400}  % data/results/golden_check/labels.csv
\newcommand{\GoldenIndianDFourRtiActRateCI}{[0.219, 0.613]}  % data/results/golden_check/labels.csv
\newcommand{\GoldenIndianDFourRtiActUnsure}{0}  % data/results/golden_check/labels.csv
\newcommand{\GoldenIndianDOneInvalid}{3}  % data/results/golden_check/labels.csv
\newcommand{\GoldenIndianDOneLabelled}{20}  % data/results/golden_check/labels.csv
\newcommand{\GoldenIndianDOneRate}{0.150}  % data/results/golden_check/labels.csv
\newcommand{\GoldenIndianDOneRateCI}{[0.052, 0.360]}  % data/results/golden_check/labels.csv
\newcommand{\GoldenIndianDOneUnsure}{8}  % data/results/golden_check/labels.csv
\newcommand{\GoldenIndianDThreeFundamentalRulesInvalid}{9}  % data/results/golden_check/labels.csv
\newcommand{\GoldenIndianDThreeFundamentalRulesLabelled}{20}  % data/results/golden_check/labels.csv
\newcommand{\GoldenIndianDThreeFundamentalRulesRate}{0.450}  % data/results/golden_check/labels.csv
\newcommand{\GoldenIndianDThreeFundamentalRulesRateCI}{[0.258, 0.658]}  % data/results/golden_check/labels.csv
\newcommand{\GoldenIndianDThreeFundamentalRulesUnsure}{0}  % data/results/golden_check/labels.csv
\newcommand{\GoldenInvalid}{20}  % data/results/golden_check/labels.csv
\newcommand{\GoldenOverallInvalid}{20}  % data/results/golden_check/labels.csv
\newcommand{\GoldenOverallLabelled}{60}  % data/results/golden_check/labels.csv
\newcommand{\GoldenOverallRate}{0.333}  % data/results/golden_check/labels.csv
\newcommand{\GoldenOverallRateCI}{[0.227, 0.459]}  % data/results/golden_check/labels.csv
\newcommand{\GoldenOverallUnsure}{8}  % data/results/golden_check/labels.csv
\newcommand{\GoldenSampled}{60}  % data/results/golden_check/labels.csv
\newcommand{\GoldenUnsure}{8}  % data/results/golden_check/labels.csv
\newcommand{\HeldOutSetups}{98}  % runs/P4-FINAL-20261106-1.json
\newcommand{\JudgeCases}{40}  % data/results/aq/judge_check.csv
\newcommand{\JudgeDisagreements}{12}  % reconstruction from judge_check.csv, judgements.jsonl, manual_facts.csv
\newcommand{\JudgeReconKappaHigh}{0.46}  % reconstruction from judge_check.csv, judgements.jsonl, manual_facts.csv
\newcommand{\JudgeReconKappaLow}{0.45}  % reconstruction from judge_check.csv, judgements.jsonl, manual_facts.csv
\newcommand{\JudgeSecondAgreement}{95}  % data/results/aq/judge_check.csv (percent)
\newcommand{\JudgeSecondKappa}{0.86}  % data/results/aq/judge_check.csv
\newcommand{\ManualJudgeFirstPassAgreement}{70}  % spec/spec_v1.2.md, AQ judge check (2026-11-12) @ a6278ca (percent)
\newcommand{\ManualJudgeFirstPassKappa}{0.48}  % spec/spec_v1.2.md, AQ judge check (2026-11-12) @ a6278ca
\newcommand{\ManualJudgeFirstPassLongCorrect}{11}  % spec/spec_v1.2.md, AQ judge check (2026-11-12) @ a6278ca
\newcommand{\ManualOcrVariantsIdenticalDocuments}{10}  % spec/spec_v1.2.md, RQ5(a) result (2026-11-10) @ 7733597
\newcommand{\ManualScannedPlacementDoclingCertain}{47}  % spec/spec_v1.2.md, Correction to RQ5(a) scanned results (2026-11-11) @ 3d5407e
\newcommand{\ManualScannedPlacementDoclingExact}{49}  % spec/spec_v1.2.md, Correction to RQ5(a) scanned results (2026-11-11) @ 3d5407e
\newcommand{\ManualScannedPlacementPdfplumber}{0}  % spec/spec_v1.2.md, Correction to RQ5(a) scanned results (2026-11-11) @ 3d5407e
\newcommand{\ManualScannedPlacementTesseractCertain}{43}  % spec/spec_v1.2.md, Correction to RQ5(a) scanned results (2026-11-11) @ 3d5407e
\newcommand{\ManualScannedPlacementTesseractExact}{41}  % spec/spec_v1.2.md, Correction to RQ5(a) scanned results (2026-11-11) @ 3d5407e
\newcommand{\PdfplumberVersion}{0.11.10}  % runs/P1-LOCAL-20261106-1.json
\newcommand{\PrefixCap}{10\%}  % src/chunk/base.py and src/chunk/chunkers.py and src/eval/stats.py and scripts/p4_analyse.py and scripts/aq_run.py and scripts/golden_check.py and scripts/p3_fixedk.py at paper-v1.2
\newcommand{\RQfiveADigitalDoclingRecall}{0.827}  % data/results/rq5/P5-LOCAL-20261110-1_digital.csv
\newcommand{\RQfiveADigitalMaxAbsDiff}{0.007}  % data/results/rq5/P5-LOCAL-20261110-1_digital.csv
\newcommand{\RQfiveADigitalMinHolm}{0.29}  % data/results/rq5/P5-LOCAL-20261110-1_digital.csv
\newcommand{\RQfiveADigitalPdfplumberRecall}{0.797--0.798}  % data/results/rq5/P5-LOCAL-20261110-1_digital.csv
\newcommand{\RQfiveADigitalQuestions}{306--308}  % data/results/rq5/P5-LOCAL-20261110-1_digital.csv
\newcommand{\RQfiveAScannedCI}{[$-$0.05, 0.08]}  % data/results/rq5/P5-LOCAL-20261111-1_scanned.csv
\newcommand{\RQfiveAScannedDiff}{+0.02}  % data/results/rq5/P5-LOCAL-20261111-1_scanned.csv
\newcommand{\RQfiveAScannedDocuments}{8}  % data/results/rq5/P5-LOCAL-20261111-1_scanned.csv
\newcommand{\RQfiveAScannedHolm}{0.50}  % data/results/rq5/P5-LOCAL-20261111-1_scanned.csv
\newcommand{\RQfiveAScannedQuestions}{34}  % data/results/rq5/P5-LOCAL-20261111-1_scanned.csv
\newcommand{\RQfiveBComparisons}{8}  % data/results/rq5/P5-LOCAL-20261110-2_chunker.csv
\newcommand{\RQfiveBFixedTokenDoclingCI}{[$-$0.027, 0.121]}  % data/results/rq5/P5-LOCAL-20261110-2_chunker.csv
\newcommand{\RQfiveBFixedTokenDoclingDefault}{0.687}  % data/results/rq5/P5-LOCAL-20261110-2_chunker.csv
\newcommand{\RQfiveBFixedTokenDoclingDiff}{+0.042}  % data/results/rq5/P5-LOCAL-20261110-2_chunker.csv
\newcommand{\RQfiveBFixedTokenDoclingHolm}{0.765}  % data/results/rq5/P5-LOCAL-20261110-2_chunker.csv
\newcommand{\RQfiveBFixedTokenDoclingTuned}{0.729}  % data/results/rq5/P5-LOCAL-20261110-2_chunker.csv
\newcommand{\RQfiveBFixedTokenPdfplumberCI}{[$-$0.001, 0.070]}  % data/results/rq5/P5-LOCAL-20261110-2_chunker.csv
\newcommand{\RQfiveBFixedTokenPdfplumberDefault}{0.707}  % data/results/rq5/P5-LOCAL-20261110-2_chunker.csv
\newcommand{\RQfiveBFixedTokenPdfplumberDiff}{+0.035}  % data/results/rq5/P5-LOCAL-20261110-2_chunker.csv
\newcommand{\RQfiveBFixedTokenPdfplumberHolm}{0.306}  % data/results/rq5/P5-LOCAL-20261110-2_chunker.csv
\newcommand{\RQfiveBFixedTokenPdfplumberTuned}{0.742}  % data/results/rq5/P5-LOCAL-20261110-2_chunker.csv
\newcommand{\RQfiveBLangchainRecursiveDoclingCI}{[0.057, 0.196]}  % data/results/rq5/P5-LOCAL-20261110-2_chunker.csv
\newcommand{\RQfiveBLangchainRecursiveDoclingDefault}{0.610}  % data/results/rq5/P5-LOCAL-20261110-2_chunker.csv
\newcommand{\RQfiveBLangchainRecursiveDoclingDiff}{+0.123}  % data/results/rq5/P5-LOCAL-20261110-2_chunker.csv
\newcommand{\RQfiveBLangchainRecursiveDoclingHolm}{$<$0.01}  % data/results/rq5/P5-LOCAL-20261110-2_chunker.csv
\newcommand{\RQfiveBLangchainRecursiveDoclingTuned}{0.733}  % data/results/rq5/P5-LOCAL-20261110-2_chunker.csv
\newcommand{\RQfiveBLangchainRecursivePdfplumberCI}{[0.038, 0.186]}  % data/results/rq5/P5-LOCAL-20261110-2_chunker.csv
\newcommand{\RQfiveBLangchainRecursivePdfplumberDefault}{0.648}  % data/results/rq5/P5-LOCAL-20261110-2_chunker.csv
\newcommand{\RQfiveBLangchainRecursivePdfplumberDiff}{+0.100}  % data/results/rq5/P5-LOCAL-20261110-2_chunker.csv
\newcommand{\RQfiveBLangchainRecursivePdfplumberHolm}{$<$0.01}  % data/results/rq5/P5-LOCAL-20261110-2_chunker.csv
\newcommand{\RQfiveBLangchainRecursivePdfplumberTuned}{0.748}  % data/results/rq5/P5-LOCAL-20261110-2_chunker.csv
\newcommand{\RQfiveBMethodOneTreeDoclingCI}{[$-$0.021, 0.053]}  % data/results/rq5/P5-LOCAL-20261110-2_chunker.csv
\newcommand{\RQfiveBMethodOneTreeDoclingDefault}{0.736}  % data/results/rq5/P5-LOCAL-20261110-2_chunker.csv
\newcommand{\RQfiveBMethodOneTreeDoclingDiff}{+0.016}  % data/results/rq5/P5-LOCAL-20261110-2_chunker.csv
\newcommand{\RQfiveBMethodOneTreeDoclingHolm}{0.816}  % data/results/rq5/P5-LOCAL-20261110-2_chunker.csv
\newcommand{\RQfiveBMethodOneTreeDoclingTuned}{0.751}  % data/results/rq5/P5-LOCAL-20261110-2_chunker.csv
\newcommand{\RQfiveBMethodOneTreePdfplumberCI}{[$-$0.000, 0.104]}  % data/results/rq5/P5-LOCAL-20261110-2_chunker.csv
\newcommand{\RQfiveBMethodOneTreePdfplumberDefault}{0.716}  % data/results/rq5/P5-LOCAL-20261110-2_chunker.csv
\newcommand{\RQfiveBMethodOneTreePdfplumberDiff}{+0.047}  % data/results/rq5/P5-LOCAL-20261110-2_chunker.csv
\newcommand{\RQfiveBMethodOneTreePdfplumberHolm}{0.306}  % data/results/rq5/P5-LOCAL-20261110-2_chunker.csv
\newcommand{\RQfiveBMethodOneTreePdfplumberTuned}{0.762}  % data/results/rq5/P5-LOCAL-20261110-2_chunker.csv
\newcommand{\RQfiveBRecursiveDoclingCI}{[$-$0.053, 0.024]}  % data/results/rq5/P5-LOCAL-20261110-2_chunker.csv
\newcommand{\RQfiveBRecursiveDoclingDefault}{0.736}  % data/results/rq5/P5-LOCAL-20261110-2_chunker.csv
\newcommand{\RQfiveBRecursiveDoclingDiff}{$-$0.015}  % data/results/rq5/P5-LOCAL-20261110-2_chunker.csv
\newcommand{\RQfiveBRecursiveDoclingHolm}{0.816}  % data/results/rq5/P5-LOCAL-20261110-2_chunker.csv
\newcommand{\RQfiveBRecursiveDoclingTuned}{0.722}  % data/results/rq5/P5-LOCAL-20261110-2_chunker.csv
\newcommand{\RQfiveBRecursivePdfplumberCI}{[$-$0.012, 0.058]}  % data/results/rq5/P5-LOCAL-20261110-2_chunker.csv
\newcommand{\RQfiveBRecursivePdfplumberDefault}{0.714}  % data/results/rq5/P5-LOCAL-20261110-2_chunker.csv
\newcommand{\RQfiveBRecursivePdfplumberDiff}{+0.025}  % data/results/rq5/P5-LOCAL-20261110-2_chunker.csv
\newcommand{\RQfiveBRecursivePdfplumberHolm}{0.760}  % data/results/rq5/P5-LOCAL-20261110-2_chunker.csv
\newcommand{\RQfiveBRecursivePdfplumberTuned}{0.738}  % data/results/rq5/P5-LOCAL-20261110-2_chunker.csv
\newcommand{\RQfourBest}{docling / method1\_tree T256}  % data/results/aq/analysis_held_out.csv
\newcommand{\RQfourBestFOne}{0.393}  % data/results/aq/analysis_held_out.csv
\newcommand{\RQfourComparisons}{6}  % data/results/aq/analysis_held_out_vs_best.csv
\newcommand{\RQfourDoclingFixedTokenSizeOneZeroTwoFourCI}{[0.050, 0.210]}  % data/results/aq/analysis_held_out_vs_best.csv
\newcommand{\RQfourDoclingFixedTokenSizeOneZeroTwoFourDiff}{+0.123}  % data/results/aq/analysis_held_out_vs_best.csv
\newcommand{\RQfourDoclingFixedTokenSizeOneZeroTwoFourHolm}{$<$0.01}  % data/results/aq/analysis_held_out_vs_best.csv
\newcommand{\RQfourDoclingFixedTokenSizeOneZeroTwoFourSetup}{docling / fixed\_token 1024}  % data/results/aq/analysis_held_out_vs_best.csv
\newcommand{\RQfourDoclingLangchainDefaultCI}{[0.002, 0.099]}  % data/results/aq/analysis_held_out_vs_best.csv
\newcommand{\RQfourDoclingLangchainDefaultDiff}{+0.048}  % data/results/aq/analysis_held_out_vs_best.csv
\newcommand{\RQfourDoclingLangchainDefaultHolm}{0.176}  % data/results/aq/analysis_held_out_vs_best.csv
\newcommand{\RQfourDoclingLangchainDefaultSetup}{docling / langchain\_default}  % data/results/aq/analysis_held_out_vs_best.csv
\newcommand{\RQfourDoclingLangchainRecursiveSizeFiveOneTwoCI}{[$-$0.033, 0.051]}  % data/results/aq/analysis_held_out_vs_best.csv
\newcommand{\RQfourDoclingLangchainRecursiveSizeFiveOneTwoDiff}{+0.010}  % data/results/aq/analysis_held_out_vs_best.csv
\newcommand{\RQfourDoclingLangchainRecursiveSizeFiveOneTwoHolm}{1.000}  % data/results/aq/analysis_held_out_vs_best.csv
\newcommand{\RQfourDoclingLangchainRecursiveSizeFiveOneTwoSetup}{docling / langchain\_recursive 512}  % data/results/aq/analysis_held_out_vs_best.csv
\newcommand{\RQfourDoclingLangchainRecursiveSizeTwoFiveSixCI}{[$-$0.021, 0.117]}  % data/results/aq/analysis_held_out_vs_best.csv
\newcommand{\RQfourDoclingLangchainRecursiveSizeTwoFiveSixDiff}{+0.035}  % data/results/aq/analysis_held_out_vs_best.csv
\newcommand{\RQfourDoclingLangchainRecursiveSizeTwoFiveSixHolm}{1.000}  % data/results/aq/analysis_held_out_vs_best.csv
\newcommand{\RQfourDoclingLangchainRecursiveSizeTwoFiveSixSetup}{docling / langchain\_recursive 256}  % data/results/aq/analysis_held_out_vs_best.csv
\newcommand{\RQfourDoclingMethodOneTreeTThreeEightFourCI}{[0.010, 0.137]}  % data/results/aq/analysis_held_out_vs_best.csv
\newcommand{\RQfourDoclingMethodOneTreeTThreeEightFourDiff}{+0.065}  % data/results/aq/analysis_held_out_vs_best.csv
\newcommand{\RQfourDoclingMethodOneTreeTThreeEightFourHolm}{0.025}  % data/results/aq/analysis_held_out_vs_best.csv
\newcommand{\RQfourDoclingMethodOneTreeTThreeEightFourSetup}{docling / method1\_tree T384}  % data/results/aq/analysis_held_out_vs_best.csv
\newcommand{\RQfourDoclingRecursiveSizeFiveOneTwoCI}{[$-$0.028, 0.113]}  % data/results/aq/analysis_held_out_vs_best.csv
\newcommand{\RQfourDoclingRecursiveSizeFiveOneTwoDiff}{+0.030}  % data/results/aq/analysis_held_out_vs_best.csv
\newcommand{\RQfourDoclingRecursiveSizeFiveOneTwoHolm}{1.000}  % data/results/aq/analysis_held_out_vs_best.csv
\newcommand{\RQfourDoclingRecursiveSizeFiveOneTwoSetup}{docling / recursive 512}  % data/results/aq/analysis_held_out_vs_best.csv
\newcommand{\RQfourScannedAboveDocling}{4}  % data/results/aq/analysis_scanned_vs_best.csv
\newcommand{\RQfourScannedAboveDoclingOcrTesseract}{5}  % data/results/aq/analysis_scanned_vs_best.csv
\newcommand{\RQfourScannedBest}{docling / recursive 256}  % data/results/aq/analysis_scanned.csv
\newcommand{\RQfourScannedBestFOne}{0.545}  % data/results/aq/analysis_scanned.csv
\newcommand{\RQfourScannedComparisons}{9}  % data/results/aq/analysis_scanned_vs_best.csv
\newcommand{\RQfourScannedCountDocling}{4}  % data/results/aq/analysis_scanned_vs_best.csv
\newcommand{\RQfourScannedCountDoclingOcrTesseract}{5}  % data/results/aq/analysis_scanned_vs_best.csv
\newcommand{\RQfourScannedSignificant}{1}  % data/results/aq/analysis_scanned_vs_best.csv
\newcommand{\RQfourScannedTopCI}{[0.065, 0.420]}  % data/results/aq/analysis_scanned_vs_best.csv
\newcommand{\RQfourScannedTopDiff}{+0.243}  % data/results/aq/analysis_scanned_vs_best.csv
\newcommand{\RQfourScannedTopHolm}{0.036}  % data/results/aq/analysis_scanned_vs_best.csv
\newcommand{\RQfourScannedTopSetup}{docling\_ocr\_tesseract / fixed\_token 256}  % data/results/aq/analysis_scanned_vs_best.csv
\newcommand{\RQfourSpearmanHeldOutFOne}{0.43}  % data/results/aq/answers_f1.csv, data/results/p3/P3-FINAL-20261106-1.csv.gz
\newcommand{\RQfourSpearmanHeldOutJudge}{0.64}  % data/results/aq/judgements.jsonl, data/results/p3/P3-FINAL-20261106-1.csv.gz
\newcommand{\RQfourSpearmanTuningFOne}{0.75}  % data/results/aq/answers_f1.csv, data/results/p3/P3-LOCAL-20261105-4.csv.gz
\newcommand{\RQfourSpearmanTuningJudge}{0.75}  % data/results/aq/judgements.jsonl, data/results/p3/P3-LOCAL-20261105-4.csv.gz
\newcommand{\RQfourTuningMinHolm}{1.00}  % data/results/aq/analysis_tuning_vs_best.csv
\newcommand{\RQoneBest}{pdfplumber / method1\_tree T384}  % runs/P4-FINAL-20261106-1.json
\newcommand{\RQoneBestRecall}{0.772}  % runs/P4-FINAL-20261106-1.json
\newcommand{\RQoneDecision}{NO-GO}  % runs/P4-FINAL-20261106-1.json
\newcommand{\RQoneDocuments}{17}  % runs/P4-FINAL-20261106-1.json
\newcommand{\RQoneFullGap}{$-$0.056}  % runs/P4-FINAL-20261106-1.json
\newcommand{\RQoneFullGapCI}{[$-$0.108, $-$0.010]}  % runs/P4-FINAL-20261106-1.json
\newcommand{\RQoneHRCheckExcludedQuestions}{7}  % runs/P4-H2-20261124-5.json
\newcommand{\RQoneHRCheckPages}{11}  % runs/P4-H2-20261124-5.json
\newcommand{\RQoneHRCheckWithDecision}{NO-GO}  % data/results/h2/P4-H2-20261124-5_exclusion.csv
\newcommand{\RQoneHRCheckWithGap}{$-$0.042}  % data/results/h2/P4-H2-20261124-5_exclusion.csv
\newcommand{\RQoneHRCheckWithGapCI}{[$-$0.093, 0.006]}  % data/results/h2/P4-H2-20261124-5_exclusion.csv
\newcommand{\RQoneHRCheckWithQuestions}{229}  % data/results/h2/P4-H2-20261124-5_exclusion.csv
\newcommand{\RQoneHRCheckWithoutDecision}{NO-GO}  % data/results/h2/P4-H2-20261124-5_exclusion.csv
\newcommand{\RQoneHRCheckWithoutGap}{$-$0.053}  % data/results/h2/P4-H2-20261124-5_exclusion.csv
\newcommand{\RQoneHRCheckWithoutGapCI}{[$-$0.105, $-$0.005]}  % data/results/h2/P4-H2-20261124-5_exclusion.csv
\newcommand{\RQoneHRCheckWithoutQuestions}{222}  % data/results/h2/P4-H2-20261124-5_exclusion.csv
\newcommand{\RQoneHTwoCertainOnlyBest}{docling / langchain\_recursive 256}  % runs/P4-H2-20261124-3.json
\newcommand{\RQoneHTwoCertainOnlyBestRecall}{0.916}  % runs/P4-H2-20261124-3.json
\newcommand{\RQoneHTwoCertainOnlyDecision}{NO-GO}  % runs/P4-H2-20261124-3.json
\newcommand{\RQoneHTwoCertainOnlyDocuments}{29}  % runs/P4-H2-20261124-3.json
\newcommand{\RQoneHTwoCertainOnlyGap}{$-$0.021}  % runs/P4-H2-20261124-3.json
\newcommand{\RQoneHTwoCertainOnlyGapCI}{[$-$0.058, 0.018]}  % runs/P4-H2-20261124-3.json
\newcommand{\RQoneHTwoCertainOnlyOracle}{0.895}  % runs/P4-H2-20261124-3.json
\newcommand{\RQoneHTwoCertainOnlyQuestions}{95}  % runs/P4-H2-20261124-3.json
\newcommand{\RQoneHTwoCertainOnlyTied}{7}  % runs/P4-H2-20261124-3.json
\newcommand{\RQoneHTwoCertainOnlyTuningWinnerRank}{1}  % data/results/p4/P4-H2-20261124-3_primary_setups.csv
\newcommand{\RQoneHTwoCertainOnlyTuningWinnerRecall}{0.916}  % data/results/p4/P4-H2-20261124-3_primary_setups.csv
\newcommand{\RQoneHTwoCertainOnlyTuningWinnerTied}{tied}  % data/results/p4/P4-H2-20261124-3_primary_setups.csv
\newcommand{\RQoneHTwoDropRuleBest}{docling / method1\_tree T256}  % runs/P4-H2-20261124-1.json
\newcommand{\RQoneHTwoDropRuleBestRecall}{0.836}  % runs/P4-H2-20261124-1.json
\newcommand{\RQoneHTwoDropRuleDecision}{NO-GO}  % runs/P4-H2-20261124-1.json
\newcommand{\RQoneHTwoDropRuleDocuments}{52}  % runs/P4-H2-20261124-1.json
\newcommand{\RQoneHTwoDropRuleGap}{$-$0.045}  % runs/P4-H2-20261124-1.json
\newcommand{\RQoneHTwoDropRuleGapCI}{[$-$0.096, 0.002]}  % runs/P4-H2-20261124-1.json
\newcommand{\RQoneHTwoDropRuleOracle}{0.792}  % runs/P4-H2-20261124-1.json
\newcommand{\RQoneHTwoDropRuleQuestions}{229}  % runs/P4-H2-20261124-1.json
\newcommand{\RQoneHTwoDropRuleTied}{16}  % runs/P4-H2-20261124-1.json
\newcommand{\RQoneHTwoDropRuleTuningWinnerRank}{2}  % data/results/p4/P4-H2-20261124-1_primary_setups.csv
\newcommand{\RQoneHTwoDropRuleTuningWinnerRecall}{0.828}  % data/results/p4/P4-H2-20261124-1_primary_setups.csv
\newcommand{\RQoneHTwoDropRuleTuningWinnerTied}{tied}  % data/results/p4/P4-H2-20261124-1_primary_setups.csv
\newcommand{\RQoneHTwoSensitivityBest}{docling / method1\_tree T256}  % runs/P4-H2-20261124-2.json
\newcommand{\RQoneHTwoSensitivityBestRecall}{0.838}  % runs/P4-H2-20261124-2.json
\newcommand{\RQoneHTwoSensitivityDecision}{NO-GO}  % runs/P4-H2-20261124-2.json
\newcommand{\RQoneHTwoSensitivityDocuments}{52}  % runs/P4-H2-20261124-2.json
\newcommand{\RQoneHTwoSensitivityGap}{$-$0.049}  % runs/P4-H2-20261124-2.json
\newcommand{\RQoneHTwoSensitivityGapCI}{[$-$0.101, $-$0.001]}  % runs/P4-H2-20261124-2.json
\newcommand{\RQoneHTwoSensitivityOracle}{0.789}  % runs/P4-H2-20261124-2.json
\newcommand{\RQoneHTwoSensitivityQuestions}{240}  % runs/P4-H2-20261124-2.json
\newcommand{\RQoneHTwoSensitivityTied}{11}  % runs/P4-H2-20261124-2.json
\newcommand{\RQoneHTwoSensitivityTuningWinnerRank}{3}  % data/results/p4/P4-H2-20261124-2_primary_setups.csv
\newcommand{\RQoneHTwoSensitivityTuningWinnerRecall}{0.826}  % data/results/p4/P4-H2-20261124-2_primary_setups.csv
\newcommand{\RQoneHTwoSensitivityTuningWinnerTied}{tied}  % data/results/p4/P4-H2-20261124-2_primary_setups.csv
\newcommand{\RQoneOracle}{0.765}  % runs/P4-FINAL-20261106-1.json
\newcommand{\RQoneOriginalGap}{$-$0.042}  % runs/P4-FINAL-20261106-1.json
\newcommand{\RQoneOriginalGapCI}{[$-$0.084, 0.002]}  % runs/P4-FINAL-20261106-1.json
\newcommand{\RQonePooledDecision}{NO-GO}  % data/results/h2/P4-H2-20261124-6_pooled.csv
\newcommand{\RQonePooledDocuments}{69}  % data/results/h2/P4-H2-20261124-6_pooled.csv
\newcommand{\RQonePooledGap}{$-$0.014}  % data/results/h2/P4-H2-20261124-6_pooled.csv
\newcommand{\RQonePooledGapCI}{[$-$0.059, 0.030]}  % data/results/h2/P4-H2-20261124-6_pooled.csv
\newcommand{\RQonePooledOracle}{0.786}  % data/results/h2/P4-H2-20261124-6_pooled.csv
\newcommand{\RQonePooledQuestions}{759}  % data/results/h2/P4-H2-20261124-6_pooled.csv
\newcommand{\RQonePooledTuningWinnerRecall}{0.800}  % data/results/h2/P4-H2-20261124-6_pooled.csv
\newcommand{\RQonePrimaryGap}{$-$0.007}  % runs/P4-FINAL-20261106-1.json
\newcommand{\RQonePrimaryGapCI}{[$-$0.045, 0.023]}  % runs/P4-FINAL-20261106-1.json
\newcommand{\RQonePrimarySentenceWindow}{5}  % data/results/p4/P4-FINAL-20261106-1_primary_setups.csv
\newcommand{\RQonePrimarySettings}{13}  % data/results/p4/P4-FINAL-20261106-1_primary_setups.csv
\newcommand{\RQonePrimarySizes}{256, 384 and 512}  % data/results/p4/P4-FINAL-20261106-1_primary_setups.csv
\newcommand{\RQoneQuestions}{530}  % runs/P4-FINAL-20261106-1.json
\newcommand{\RQoneRawGap}{+0.058}  % runs/P4-FINAL-20261106-1.json
\newcommand{\RQoneSetups}{26}  % runs/P4-FINAL-20261106-1.json
\newcommand{\RQoneTied}{15}  % runs/P4-FINAL-20261106-1.json
\newcommand{\RQoneTuningWinner}{docling / langchain\_recursive 256}  % data/results/p4/P4-NOINDIAN-20261114-1_p4.csv
\newcommand{\RQoneTuningWinnerRank}{20}  % data/results/p4/P4-FINAL-20261106-1_primary_setups.csv
\newcommand{\RQoneTuningWinnerRecall}{0.715}  % data/results/p4/P4-FINAL-20261106-1_primary_setups.csv
\newcommand{\RQoneTuningWinnerTied}{tied}  % data/results/p4/P4-FINAL-20261106-1_primary_setups.csv
\newcommand{\RQthreeDoclingHitOneBest}{method1\_tree T384}  % data/results/p3/fixedk_tuning.csv
\newcommand{\RQthreeDoclingHitOneScore}{0.622}  % data/results/p3/fixedk_tuning.csv
\newcommand{\RQthreeDoclingHitOneSize}{295}  % data/results/p3/fixedk_tuning.csv
\newcommand{\RQthreeDoclingMrrBest}{method1\_tree T1024}  % data/results/p3/fixedk_tuning.csv
\newcommand{\RQthreeDoclingMrrScore}{0.742}  % data/results/p3/fixedk_tuning.csv
\newcommand{\RQthreeDoclingMrrSize}{715}  % data/results/p3/fixedk_tuning.csv
\newcommand{\RQthreeDoclingRecallFiveBest}{fixed\_token 1024}  % data/results/p3/fixedk_tuning.csv
\newcommand{\RQthreeDoclingRecallFiveScore}{0.904}  % data/results/p3/fixedk_tuning.csv
\newcommand{\RQthreeDoclingRecallFiveSize}{978}  % data/results/p3/fixedk_tuning.csv
\newcommand{\RQthreeDoclingRecallFourZeroZeroZeroBest}{method1\_tree T384 prefix}  % data/results/p3/fixedk_tuning.csv
\newcommand{\RQthreeDoclingRecallFourZeroZeroZeroScore}{0.925}  % data/results/p3/fixedk_tuning.csv
\newcommand{\RQthreeDoclingRecallFourZeroZeroZeroSize}{309}  % data/results/p3/fixedk_tuning.csv
\newcommand{\RQthreeDoclingRecallOneZeroBest}{fixed\_token 896}  % data/results/p3/fixedk_tuning.csv
\newcommand{\RQthreeDoclingRecallOneZeroScore}{0.964}  % data/results/p3/fixedk_tuning.csv
\newcommand{\RQthreeDoclingRecallOneZeroSize}{879}  % data/results/p3/fixedk_tuning.csv
\newcommand{\RQthreePdfplumberHitOneBest}{recursive 1024}  % data/results/p3/fixedk_tuning.csv
\newcommand{\RQthreePdfplumberHitOneScore}{0.632}  % data/results/p3/fixedk_tuning.csv
\newcommand{\RQthreePdfplumberHitOneSize}{820}  % data/results/p3/fixedk_tuning.csv
\newcommand{\RQthreePdfplumberMrrBest}{recursive 1024}  % data/results/p3/fixedk_tuning.csv
\newcommand{\RQthreePdfplumberMrrScore}{0.748}  % data/results/p3/fixedk_tuning.csv
\newcommand{\RQthreePdfplumberMrrSize}{820}  % data/results/p3/fixedk_tuning.csv
\newcommand{\RQthreePdfplumberRecallFiveBest}{fixed\_token 1024}  % data/results/p3/fixedk_tuning.csv
\newcommand{\RQthreePdfplumberRecallFiveScore}{0.904}  % data/results/p3/fixedk_tuning.csv
\newcommand{\RQthreePdfplumberRecallFiveSize}{980}  % data/results/p3/fixedk_tuning.csv
\newcommand{\RQthreePdfplumberRecallFourZeroZeroZeroBest}{sentence 5}  % data/results/p3/fixedk_tuning.csv
\newcommand{\RQthreePdfplumberRecallFourZeroZeroZeroScore}{0.928}  % data/results/p3/fixedk_tuning.csv
\newcommand{\RQthreePdfplumberRecallFourZeroZeroZeroSize}{249}  % data/results/p3/fixedk_tuning.csv
\newcommand{\RQthreePdfplumberRecallOneZeroBest}{fixed\_token 896}  % data/results/p3/fixedk_tuning.csv
\newcommand{\RQthreePdfplumberRecallOneZeroScore}{0.968}  % data/results/p3/fixedk_tuning.csv
\newcommand{\RQthreePdfplumberRecallOneZeroSize}{875}  % data/results/p3/fixedk_tuning.csv
\newcommand{\RQthreeSettingsKept}{50}  % data/results/p3/fixedk_tuning.csv
\newcommand{\RQthreeSettingsLeftOut}{13}  % data/results/p3/fixedk_tuning.csv
\newcommand{\RQtwoBestSizes}{209--439}  % data/results/p3/size_by_source_summary_P3-LOCAL-20261105-4.csv
\newcommand{\RQtwoHTwoCertainOnlyDoclingOneTwoFiveZero}{$-$0.582}  % data/results/h2/P4-H2-20261124-7_rq2.csv
\newcommand{\RQtwoHTwoCertainOnlyDoclingSevenFiveZero}{$-$0.813}  % data/results/h2/P4-H2-20261124-7_rq2.csv
\newcommand{\RQtwoHTwoCertainOnlyDoclingTwoZeroZeroZero}{$-$0.632}  % data/results/h2/P4-H2-20261124-7_rq2.csv
\newcommand{\RQtwoHTwoCertainOnlyPdfplumberOneTwoFiveZero}{$-$0.484}  % data/results/h2/P4-H2-20261124-7_rq2.csv
\newcommand{\RQtwoHTwoCertainOnlyPdfplumberSevenFiveZero}{$-$0.681}  % data/results/h2/P4-H2-20261124-7_rq2.csv
\newcommand{\RQtwoHTwoCertainOnlyPdfplumberTwoZeroZeroZero}{$-$0.055}  % data/results/h2/P4-H2-20261124-7_rq2.csv
\newcommand{\RQtwoHTwoDropRuleDoclingOneTwoFiveZero}{$-$0.830}  % data/results/h2/P4-H2-20261124-7_rq2.csv
\newcommand{\RQtwoHTwoDropRuleDoclingSevenFiveZero}{$-$0.885}  % data/results/h2/P4-H2-20261124-7_rq2.csv
\newcommand{\RQtwoHTwoDropRuleDoclingTwoZeroZeroZero}{$-$0.148}  % data/results/h2/P4-H2-20261124-7_rq2.csv
\newcommand{\RQtwoHTwoDropRulePdfplumberOneTwoFiveZero}{$-$0.599}  % data/results/h2/P4-H2-20261124-7_rq2.csv
\newcommand{\RQtwoHTwoDropRulePdfplumberSevenFiveZero}{$-$0.714}  % data/results/h2/P4-H2-20261124-7_rq2.csv
\newcommand{\RQtwoHTwoDropRulePdfplumberTwoZeroZeroZero}{$-$0.231}  % data/results/h2/P4-H2-20261124-7_rq2.csv
\newcommand{\RQtwoHTwoLargestBudgetIncludeZero}{5}  % data/results/h2/P4-H2-20261124-7_rq2.csv
\newcommand{\RQtwoHTwoLargestBudgetValues}{6}  % data/results/h2/P4-H2-20261124-7_rq2.csv
\newcommand{\RQtwoHTwoNegative}{18}  % data/results/h2/P4-H2-20261124-7_rq2.csv
\newcommand{\RQtwoHTwoSensitivityDoclingOneTwoFiveZero}{$-$0.764}  % data/results/h2/P4-H2-20261124-7_rq2.csv
\newcommand{\RQtwoHTwoSensitivityDoclingSevenFiveZero}{$-$0.885}  % data/results/h2/P4-H2-20261124-7_rq2.csv
\newcommand{\RQtwoHTwoSensitivityDoclingTwoZeroZeroZero}{$-$0.352}  % data/results/h2/P4-H2-20261124-7_rq2.csv
\newcommand{\RQtwoHTwoSensitivityPdfplumberOneTwoFiveZero}{$-$0.516}  % data/results/h2/P4-H2-20261124-7_rq2.csv
\newcommand{\RQtwoHTwoSensitivityPdfplumberSevenFiveZero}{$-$0.665}  % data/results/h2/P4-H2-20261124-7_rq2.csv
\newcommand{\RQtwoHTwoSensitivityPdfplumberTwoZeroZeroZero}{$-$0.225}  % data/results/h2/P4-H2-20261124-7_rq2.csv
\newcommand{\RQtwoHTwoSetupsPerParser}{13}  % data/results/h2/P4-H2-20261124-7_rq2.csv
\newcommand{\RQtwoHTwoValues}{18}  % data/results/h2/P4-H2-20261124-7_rq2.csv
\newcommand{\RQtwoHeldOutOHRBestSizes}{216--318}  % data/results/p3/size_by_source_summary_P3-FINAL-20261106-1.csv
\newcommand{\RQtwoHeldOutOHRRhoRange}{$-$0.84 to $-$0.93}  % data/results/p3/size_by_source_summary_P3-FINAL-20261106-1.csv
\newcommand{\RQtwoRhoRange}{$-$0.61 to $-$0.97}  % data/results/p3/size_by_source_summary_P3-LOCAL-20261105-4.csv
\newcommand{\RRFk}{60}  % runs/P3-LOCAL-20261105-4.json
\newcommand{\Resamples}{2{,}000}  % runs/P4-FINAL-20261106-1.json
\newcommand{\RobBMTwoFiveOnlyHeldOutOneGap}{$-$0.029}  % data/results/supplementary-20261114/retriever_summary.csv
\newcommand{\RobBMTwoFiveOnlyHeldOutOneGapCI}{[$-$0.079, 0.013]}  % data/results/supplementary-20261114/retriever_summary.csv
\newcommand{\RobDenseOnlyHeldOutOneGap}{$-$0.003}  % data/results/supplementary-20261114/retriever_summary.csv
\newcommand{\RobDenseOnlyHeldOutOneGapCI}{[$-$0.074, 0.055]}  % data/results/supplementary-20261114/retriever_summary.csv
\newcommand{\RobGridSixChunkersThreeTwoTuningGap}{+0.001}  % data/results/ext/P4-EXT-20261119-2_rq1.csv
\newcommand{\RobGridSixChunkersThreeTwoTuningGapCI}{[$-$0.039, 0.042]}  % data/results/ext/P4-EXT-20261119-2_rq1.csv
\newcommand{\RobNomicDenseTuningGap}{$-$0.012}  % data/results/rob/P4-ROB-20261120-1_rq1.csv
\newcommand{\RobNomicDenseTuningGapCI}{[$-$0.056, 0.036]}  % data/results/rob/P4-ROB-20261120-1_rq1.csv
\newcommand{\RobNomicHybridTuningGap}{$-$0.000}  % data/results/rob/P4-ROB-20261120-1_rq1.csv
\newcommand{\RobNomicHybridTuningGapCI}{[$-$0.036, 0.034]}  % data/results/rob/P4-ROB-20261120-1_rq1.csv
\newcommand{\RobNomicHybridWOIndianTuningGap}{$-$0.020}  % data/results/rob/P4-ROB-20261120-1_rq1.csv
\newcommand{\RobNomicHybridWOIndianTuningGapCI}{[$-$0.056, 0.017]}  % data/results/rob/P4-ROB-20261120-1_rq1.csv
\newcommand{\RobWithoutIndianHeldOutOneGap}{$-$0.007}  % data/results/p4/P4-FINAL-NOINDIAN-20261114-1_p4.csv
\newcommand{\RobWithoutIndianHeldOutOneGapCI}{[$-$0.047, 0.028]}  % data/results/p4/P4-FINAL-NOINDIAN-20261114-1_p4.csv
\newcommand{\SanityDamageCases}{3}  % runs/tests/test_p3_sanity_floor_damage.txt
\newcommand{\SanityDamageDocuments}{11}  % runs/tests/test_p3_sanity_floor_damage.txt
\newcommand{\SanityDamageDrop}{0.08--0.10}  % runs/tests/test_p3_sanity_floor_damage.txt
\newcommand{\SanityDamageOtherMax}{0.01}  % runs/tests/test_p3_sanity_floor_damage.txt
\newcommand{\SanityDamagePassed}{3}  % runs/tests/test_p3_sanity_floor_damage.txt
\newcommand{\SanityDamageQuestions}{112}  % runs/tests/test_p3_sanity_floor_damage.txt
\newcommand{\SanityFloorCases}{6}  % runs/tests/test_p3_sanity_floor_damage.txt
\newcommand{\SanityFloorFailedOriginal}{6}  % runs/tests/test_p3_sanity_floor_damage.txt
\newcommand{\SanityFloorLower}{0.14--0.17}  % runs/tests/test_p3_sanity_floor_damage.txt
\newcommand{\SizeMax}{1{,}024}  % data/results/p4/P4-LOCAL-20261105-1_full_setups.csv
\newcommand{\SizeMin}{256}  % data/results/p4/P4-LOCAL-20261105-1_full_setups.csv; sizes shared by fixed_token, recursive, langchain_recursive, method1_tree; also {'recursive': [600]}
\newcommand{\SizeStep}{128}  % data/results/p4/P4-LOCAL-20261105-1_full_setups.csv
\newcommand{\SplitHalfReps}{200}  % runs/P4-FINAL-20261106-1.json
\newcommand{\TopKLarge}{10}  % runs/P3-FIXEDK-20261112-1.json
\newcommand{\TopKSmall}{5}  % runs/P3-FIXEDK-20261112-1.json
\newcommand{\TuningDocuments}{25}  % pdfs/split.csv
\newcommand{\TuningMainDocuments}{19}  % runs/P4-LOCAL-20261105-1.json
\newcommand{\TuningMainQuestions}{860}  % runs/P4-LOCAL-20261105-1.json
\newcommand{\TuningScoredDocuments}{20}  % data/results/p3/P3-LOCAL-20261105-4.csv.gz (hybrid, Recall@1250)
\newcommand{\TuningSetups}{100}  % runs/P4-LOCAL-20261105-1.json

% The generated figures keep their data; only their size and panel layout are set here.
\makeatletter
\newcommand{\FigHeight}{0.38\columnwidth}
\newcommand{\CompactFigure}[2]{% panel width, horizontal separation
  \pgfplotsset{
    height/.code={\pgfkeyssetvalue{/pgfplots/height}{\FigHeight}},
    width/.code={\pgfkeyssetvalue{/pgfplots/width}{#1}},
    group/group size/.code={\def\pgfplots@group@columns{2}\def\pgfplots@group@rows{1}},
    group/horizontal sep/.initial=#2,
    every axis/.append style={label style={font=\scriptsize}, tick label style={font=\scriptsize},
      title style={font=\scriptsize}, legend style={font=\scriptsize}}}}
\newcommand{\sizefigure}{\begingroup\CompactFigure{0.5\columnwidth-1.3cm}{1.1cm}
\pgfplotsset{legend style/.code={\pgfplotsset{every axis legend/.append style={##1, at={(-0.25,-0.38)}}}}}
% Generated by paper_ieee/scripts/make_artifacts.py; do not edit.
\begin{figure}[!t]
\centering
\input{generated/figures/rq2.tex}
\caption{RQ2 on tuning: Recall@1,250 against mean chunk size, per source.}
\label{fig:rq2}
\end{figure}

\endgroup}
\newcommand{\tuningfigure}{\begingroup\CompactFigure{0.5\columnwidth-2.4cm}{1.3cm}
\pgfplotsset{every axis/.append style={xlabel style={align=center, text width=2.6cm}, x tick label style={rotate=60, anchor=east}}}
% Generated by paper_ieee/scripts/make_artifacts.py; do not edit.
\begin{figure}[!t]
\centering
\input{generated/figures/rq5.tex}
\caption{RQ5: (a) parser variants minus default, digital PDFs; (b) tuned size minus default, held-out 1. Filled marks: significant after Holm.}
\label{fig:rq5}
\end{figure}

\endgroup}
\makeatother
\begin{document}
\makeatletter\def\input@path{{../paper_ieee/}}\makeatother % the generated floats input generated/figures/*

\title{What Matters When Chunking PDFs for Retrieval:\\Size, Budget, and Parser}
\author{\IEEEauthorblockN{Aman Chaurasia}
\IEEEauthorblockA{Independent Researcher\\
New Delhi, India\\
aman.1si17cs008@gmail.com}}
\maketitle
\bstctlcite{conf:control}

\begin{abstract}
Retrieval-augmented generation over PDFs depends on how each document is
parsed and cut into chunks, and a common proposal is to choose the parser and
chunker per document. We test whether this choice pays off; the
design, scorer and decision rule of this test were fixed before the first
held-out set was scored.
We score \RQoneSetups{} parse-and-chunk setups (two parsers,
\RQonePrimarySettings{} chunker settings) on two sealed held-out sets,
\RQoneDocuments{} PDFs with \RQoneQuestions{} questions with located answer evidence and
\RQoneHTwoDropRuleDocuments{} ViDoRe V3 PDFs with
\RQoneHTwoDropRuleQuestions{} questions with human-drawn evidence boxes, reading ranked chunks up to
a fixed budget of 1,250 tokens. A per-document oracle, cross-fitted on half of
each document's questions, falls below the best single setup on both sets
(gap \RQonePrimaryGap{} \RQonePrimaryGapCI{} and \RQoneHTwoDropRuleGap{}
\RQoneHTwoDropRuleGapCI{}), and the pre-registered rule gives NO-GO (no per-document chooser) on both.
We find no evidence of a gain from choosing a setup per document, although
few questions per document limit what the oracle can learn, so the gap is not
a bound on every chooser. At a budget of
750 or 1,250 tokens, smaller chunks tend to retrieve more; on tuning data,
a fixed number of chunks reverses this ranking, as large chunks put more text
in front of the reader. In a post-hoc comparison on the first held-out set,
tuned sizes beat LangChain's default chunk size by \RQfiveBLangchainRecursivePdfplumberDiff{}
to \RQfiveBLangchainRecursiveDoclingDiff{} in recall, while parser settings
show no detectable gain. We release the code, the checksummed results and a tool
with fixed defaults.
\end{abstract}

\begin{IEEEkeywords}
retrieval-augmented generation, document chunking, PDF parsing, evaluation, pre-registration
\end{IEEEkeywords}

\section{Introduction}\label{sec:intro}

Systems that answer questions over PDFs first extract the text, then cut it
into chunks that are indexed and retrieved. Because documents differ in layout
and structure, a natural idea is to choose the parser and chunker for each
document, and recent work proposes doing so from intrinsic quality
signals~\cite{demourajunior2026adaptive}. Before building such a chooser, we ask
how much retrieval quality is lost by using one fixed parse-and-chunk setup for
every PDF, and whether a choice per document, made by a cross-fitted oracle from
a part of each document's own questions, recovers any of it. The oracle selects
on few questions, so its gap is not a bound on every chooser. We study five
questions. \textbf{RQ1:} Does choosing a setup per document improve retrieval over
the best single setup? \textbf{RQ2:} At a fixed token budget, how does chunk size
relate to recall? \textbf{RQ3:} Does the best chunk size change when a fixed number
of chunks replaces the fixed token budget? \textbf{RQ4:} Do differences in answer
quality follow differences in retrieval? \textbf{RQ5:} Do tuned parser settings and
chunk sizes beat their defaults (a library default for the parsers and
LangChain, our own starting size for the other chunkers)?

Every setup reads its ranked chunks up to the same token budget. The RQ1 design,
the scorer and the RQ1 decision rule were written into dated specifications
before the first held-out set was scored; the designs of the other parts came
later and are dated in Section~\ref{sec:method}. We find no evidence of a gain from
choosing a setup per document, and the NO-GO decision is repeated on a second
held-out set with a different source of gold evidence. 


\section{Related Work}\label{sec:related}

\textbf{Chunking strategies and their selection.} Shaukat et
al.~\cite{shaukat2026systematic} test thirty-six segmentation methods with five
embedders and report that methods aware of content beat plain fixed-length
splitting. Qu et
al.~\cite{qu2025semantic} find that, over three retrieval tasks, the gains from
semantic chunking were not steady enough to pay for its computation. Other work
proposes new cutters: HiChunk~\cite{lu2025hichunk} uses fine-tuned language
models to build a hierarchy over the document, and MoC~\cite{zhao2025moc} guides
chunkers to output regular expressions . The work closest to our main question is Adaptive
Chunking~\cite{demourajunior2026adaptive}, which picks a chunker for each document
using five measures computed from the document alone, and reports better answers
than with a single method; its pool contains the LLM-regex and recursive
splitters we reuse.

\textbf{Chunk size and budgets.} 
Chen et al.~\cite{chen2024densex} index propositions and find better retrieval than with passages. Bhat et
al.~\cite{bhat2025chunksize} find that the best fixed size depends on the data and that
it also differs between embedding models. Our RQ3 offers one possible reason for such differences, which we do
not test: in our tuning data, the preferred size changes when the number of chunks, rather than
the number of tokens, is held fixed. 

\textbf{Parsing, OCR and benchmarks.} OHR-Bench~\cite{zhang2024ocr} follows OCR
noise into RAG and concludes that none of the OCR solutions it tests is good
enough for building knowledge bases. Santos et al.~\cite{santos2026pdf} compare
four conversion tools  and attribute more of
the accuracy to metadata and hierarchy-aware splitting than to the tool. El Bachyr
et al.~\cite{elbachyr2026empirical} examine parsers and chunkers jointly , and Singh~\cite{singh2026parser} runs a factorial study on
Indian regulatory texts in which parser and chunker interact.
UniDoc-Bench~\cite{peng2025unidoc} and ViDoRe~V3~\cite{loison2026vidore} supply
questions over visually rich pages; the latter also has human-drawn evidence
boxes, which we use as gold.

\textbf{This study.} We compare setups under a fixed token budget, and evaluate
the choice of a setup per document on two sealed held-out sets with a
cross-fitted oracle, against a decision rule written in advance. Studies with
other data, measures or selection signals are not directly comparable.


\section{Method}\label{sec:method}

\subsection{Pre-registration}\label{sec:prereg}

We fixed the RQ1 design, the scorer and the RQ1 decision rule before held-out
set~1 was scored; the other designs and RQ2's confirmatory rule came later. One exposure came earlier: while designing M1 we viewed the
parsed tree structure of most documents in the first part of held-out set~1 (no
questions, answers or scores). Held-out set~1 was scored once, on 6~November,
and its results were reused (RQ5(b), answer quality, further reports and
checks); no chunker in the choice sets, scorer or decision rule was changed after
it was opened. A second held-out set, added later as a replication, was scored
once per gold variant, on 24~November 2026. The timeline, all deviations and every reuse are in
the repository (tag paper-v1.4).

\subsection{Documents and questions}\label{sec:docs}

\textbf{Tuning.} \TuningDocuments{} documents, of which \TuningMainDocuments{}
(\TuningMainQuestions{} questions) enter the main analysis, from
CUAD~\cite{hendrycks2021cuad}, OHR-Bench~\cite{zhang2024ocr}, an Indian
regulatory benchmark~\cite{singh2026parser}, QASPER~\cite{dasigi2021qasper} and
UniDoc-Bench~\cite{peng2025unidoc} (the last two left out).

\textbf{Held-out set 1.} \AccountingHeldOut{} sealed documents from known
sources, new document types and out of domain (one from
FinanceBench~\cite{islam2023financebench}). \AccountingMain{} documents with
\RQoneQuestions{} questions can be scored.

\textbf{Held-out set 2.} ViDoRe~V3~\cite{loison2026vidore}, a source new to the study,
selected from metadata alone (two earlier designs were replaced before any of their
documents was drawn or opened), keeping queries whose evidence pages are all text
or table pages: \HTwoQueries{} queries in \HTwoSelectedDocuments{} documents before
gold filtering. Under the primary gold rule, \RQoneHTwoDropRuleQuestions{}
questions in \RQoneHTwoDropRuleDocuments{} documents can be scored. 

\textbf{Question quality.} The OHR-Bench, UniDoc-Bench and Indian questions are LLM-generated. The author's check (one rater, stopped after its
first batch against its own expansion rule) found \GoldenOverallInvalid{} of
\GoldenOverallLabelled{} Indian tuning questions invalid.

\subsection{Gold}\label{sec:gold}

One matching rule places gold evidence as character spans in each parser's own
output. For held-out set~1 it is each source's answer or evidence span; a question that cannot be placed with certainty is dropped. This removes
\CoverageHeldOutDoclingShare{} of held-out set~1 questions for Docling and
\CoverageHeldOutPdfplumberShare{} for pdfplumber. For held-out set~2 the evidence
is the text inside the human-drawn boxes; the primary rule keeps a
page when the matcher's frozen drop rule accepts all its box text; two
sensitivity variants keep a query if any evidence page is placed  or
apply the held-out set~1 rule (certain-only). On 23~November, after seeing how many questions each rule keeps but before any
held-out set~2 score, we made the drop rule primary. Comparisons use only questions kept for every compared setup, which hides
coverage differences between parsers.

\subsection{Setups}\label{sec:setups}

A setup pairs a parser with a chunker setting. The parsers are
Docling~2.143.0~\cite{auer2024docling} (D, default settings) and
pdfplumber~0.11.10~\cite{pdfplumber} (P). The chunkers are fixed-token (FT),
recursive (Rec), LangChain recursive~\cite{langchain} (LC), sentence groups (Sent;
Stanza~\cite{stanza}), semantic (Qwen3-Embedding-0.6B~\cite{qwen3embedding};
tuning only, out of GPU memory on held-out data), LLM-regex
(Gemini~3.5 Flash-Lite~\cite{gemini}), page-based, and our structure-tree method, M1; four reuse the code of Adaptive
Chunking~\cite{demourajunior2026adaptive}. In names, the number is the size; +p marks a prefix,
LC-def LangChain's default and D-Tess Docling with Tesseract OCR. Chunks under 100 tokens are merged and those over 1,024 tokens split again. M1 keeps a parser-tree node whole if it fits the target
size, opens it otherwise, and packs neighbouring pieces.

The primary choice set was fixed before held-out set~1 was opened: Rec, LC, FT and M1 without prefix, each at 256, 384 and
512 tokens, plus sentence groups of five. That gives \RQonePrimarySettings{}
settings and, with two parsers, \RQoneSetups{} setups. Further choice sets are the original grid (\RQoneOriginalSetups{} setups on held-out set~1) and
the full grid (\HeldOutSetups{}).

\subsection{Retrieval and scoring}\label{sec:scoring}

Retrieval is BM25~\cite{bm25}, dense (bge-m3~\cite{chen2024m3}) or a hybrid by
reciprocal rank fusion~\cite{rrf}.
Recall@$B$ is the share of gold characters, each counted once, in the ranked
chunks read until $B$ tokens are used, the last one cut at the budget. We use $B \in \{750, 1{,}250, 2{,}000\}$; the headline result
is hybrid Recall@1,250. A floor test (chunks cut at random positions score \SanityFloorLower{} lower) failed
in all \SanityFloorFailedOriginal{} parser and retriever cases under its
pre-written rule and passed only under a rule amended after that result.

RQ3 (Recall@$k$, Hit@1 and MRR at fixed $k$ on tuning data) was designed on
12~November, after the held-out set~1 results, and has no test. A later control, also designed after
the held-out set~1 results, adds Lift@5.

\subsection{Statistics}\label{sec:stats}

Question scores are averaged per document, with 95\% intervals from a document
bootstrap~\cite{bootstrap} (2,000 document resamples, seed 42, two-sided). Two setups are tied when the paired interval of their per-document
difference contains zero. The tie rule changed twice before held-out set~1 was scored.

With split-half cross-fitting (the idea of~\cite{crossfitting}, not their
estimator) over 200 random splits, the oracle picks the best setup on half of a document's questions and
is measured on the other. The gap is the mean per-document recall of the oracle minus the best single setup. A document
with one question takes the best single setup's value. This holds for
\RQoneOneQuestionDocuments{} of \RQoneDocuments{} documents on held-out set~1 and
for \RQoneHTwoOneQuestionDocuments{} of \RQoneHTwoDropRuleDocuments{} on held-out
set~2, where the median document has \RQoneHTwoMedianQuestions{} questions. Choosing the best single setup on the questions it is scored on, and fitting the
oracle on half of a document's questions, both push the gap down, so it is not an
upper bound on what any chooser could gain. The rule: a per-document chooser is worth building only if the gap's interval lies
above zero. Multiple comparisons use Holm's correction~\cite{holm}.

RQ2 began as a descriptive tuning analysis, added on 12~November after the
tuning scores were known. A confirmatory rule was fixed before held-out set~2
was scored: RQ2 holds if, for each parser and budget, Spearman's
$\rho$~\cite{spearman} of measured chunk size against recall over that parser's
\RQtwoHTwoSetupsPerParser{} primary setups is negative. These setups span target
sizes 256 to 512 only.

\subsection{RQ4 and RQ5}\label{sec:aq}

RQ5(a), designed after held-out set~1 was scored, compares parser settings with defaults on digital tuning PDFs and scanned PDFs from
MP-DocVQA~\cite{tito2023mpdocvqa}. RQ5(b), whether a tuned chunk size beats the default, was specified on
10~November after the held-out set~1 scores existed: a second use of that set, not
confirmatory. 

For RQ4, the design, fixed after held-out set~1 was scored, had one
measure, an LLM judge; token F1 and
containment, added after the answers existed, are primary, and the judge (Qwen3-8B~\cite{qwen3}) failed its
human check and is secondary only. We sample \AQSampleShare{} of each PDF's questions and compare seven Docling
setups, and five scanned-PDF settings with Docling's default OCR or
Tesseract~\cite{tesseract}. Gemini Flash-Lite answers from the 1,250 tokens
hybrid retrieval reads.


\section{Results}\label{sec:results}

\subsection{RQ1: one fixed setup versus a choice per document}\label{sec:rq1}

% Generated by paper_ieee/scripts/make_artifacts.py; do not edit.
\begin{table*}[!t]
\centering
\caption{RQ1: hybrid Recall@1,250, primary choice set unless named. Held-out set~2 rows give documents / questions. $^\dagger$Secondary: oracle minus the tuning winner, not the per-set gap; bootstrap stratified by set.}
\label{tab:rq1}
\footnotesize
\begin{tabular}{lrlrrrll}
\toprule
Choice set & Setups & Best single & Recall & Tied & Oracle & Gap [95\%] & Rule \\
\midrule
\multicolumn{8}{l}{\emph{Held-out set~1: 17 documents, 530 questions}} \\
Primary & 26 & P/M1-384 & 0.772 & 15 & 0.765 & $-$0.007 [$-$0.045, 0.023] & NO-GO \\
Full & 98 & P/M1-256+p & 0.773 & 23 & 0.718 & $-$0.056 [$-$0.108, $-$0.010] & NO-GO \\
Original & 42 & D/Sent-5 & 0.756 & 17 & 0.714 & $-$0.042 [$-$0.084, 0.002] & NO-GO \\
\midrule
\multicolumn{8}{l}{\emph{Held-out set~2 (ViDoRe V3), primary choice set}} \\
Drop rule, 52 / 229 & 26 & D/M1-256 & 0.836 & 16 & 0.792 & $-$0.045 [$-$0.096, 0.002] & NO-GO \\
Sensitivity, 52 / 240 & 26 & D/M1-256 & 0.838 & 11 & 0.789 & $-$0.049 [$-$0.101, $-$0.001] & NO-GO \\
Certain-only, 29 / 95 & 26 & D/LC-256 & 0.916 & 7 & 0.895 & $-$0.021 [$-$0.058, 0.018] & NO-GO \\
\midrule
\multicolumn{8}{l}{\emph{Pooled, secondary}} \\
Pooled$^\dagger$, 69 / 759 & 26 & D/LC-256 (fixed) & 0.800 & -- & 0.786 & $-$0.014 [$-$0.059, 0.030] & NO-GO$^\dagger$ \\
\bottomrule
\end{tabular}
\end{table*}


On held-out set~1 (\RQoneDocuments{} documents, \RQoneQuestions{} questions;
Table~\ref{tab:rq1}), the best single setup is \RQoneBest{}, with hybrid
Recall@1,250 of \RQoneBestRecall{}, and \RQoneTied{} of the \RQoneSetups{} setups,
itself included, form the statistically tied group. The cross-fitted oracle
reaches \RQoneOracle{}, a gap of \RQonePrimaryGap{} \RQonePrimaryGapCI{}. The
pre-registered rule gives \RQoneDecision{}. An oracle that picks and scores on
the same questions would show a gain of \RQoneRawGap{}. Larger choice sets do not
change the verdict (Table~\ref{tab:rq1}).

On held-out set~2, under the primary gold rule (\RQoneHTwoDropRuleDocuments{}
documents, \RQoneHTwoDropRuleQuestions{} questions), the best single setup is
\RQoneHTwoDropRuleBest{} (\RQoneHTwoDropRuleBestRecall{}), the oracle reaches
\RQoneHTwoDropRuleOracle{}, and the gap is \RQoneHTwoDropRuleGap{}
\RQoneHTwoDropRuleGapCI{}. All three gold variants are NO-GO. The primary interval
is wider than on held-out set~1 (\RQoneHTwoDropRuleGapWidth{} against
\RQonePrimaryGapWidth{}), so the second set did not narrow the estimate. A
pre-stated check without the questions on the \RQoneHRCheckPages{} HR pages whose
image and PDF shapes differ by more than 5\% is also NO-GO.

The setup chosen on tuning data, \RQoneTuningWinner{}, applied unchanged, is
statistically tied with the best single setup on every held-out set and gold
variant. Pooled over both sets (\RQonePooledDocuments{} documents,
\RQonePooledQuestions{} questions; secondary), the oracle
(\RQonePooledOracle{}) stays below the tuning winner
(\RQonePooledTuningWinnerRecall{}), a gap of \RQonePooledGap{}
\RQonePooledGapCI{}; this is a different quantity, and pooling did not narrow the
interval. 

\subsection{RQ2: chunk size at a fixed token budget}\label{sec:rq2}

% Generated by paper_ieee/scripts/make_artifacts.py; do not edit.
\begin{table*}[!t]
\centering
\caption{RQ2 on held-out set~2: Spearman's $\rho$ of chunk size against Recall@$B$ per parser, with 95\% document bootstrap intervals.}
\label{tab:rq2h2}
\footnotesize
\begin{tabular}{lllll}
\toprule
Gold & Parser & 750 & 1,250 & 2,000 \\
\midrule
Drop rule & D & $-$0.885 [$-$0.929, $-$0.313] & $-$0.830 [$-$0.901, $-$0.412] & $-$0.148 [$-$0.648, 0.495] \\
 & P & $-$0.714 [$-$0.863, $-$0.121] & $-$0.599 [$-$0.841, 0.033] & $-$0.231 [$-$0.731, 0.533] \\
Sensitivity & D & $-$0.885 [$-$0.929, $-$0.275] & $-$0.764 [$-$0.874, $-$0.368] & $-$0.352 [$-$0.692, 0.258] \\
 & P & $-$0.665 [$-$0.863, $-$0.176] & $-$0.516 [$-$0.846, 0.050] & $-$0.225 [$-$0.747, 0.528] \\
Certain-only & D & $-$0.813 [$-$0.885, $-$0.159] & $-$0.582 [$-$0.736, $-$0.236] & $-$0.632 [$-$0.681, $-$0.104] \\
 & P & $-$0.681 [$-$0.857, 0.368] & $-$0.484 [$-$0.797, 0.423] & $-$0.055 [$-$0.720, 0.506] \\
\bottomrule
\end{tabular}
\end{table*}

\sizefigure

At a fixed budget of 750 or 1,250 tokens, smaller chunks tend to retrieve more of
the gold text; at 2,000 tokens this is not established on held-out set~2. On
tuning data, the descriptive analysis finds Spearman's $\rho$ between measured
chunk size and recall negative for every parser, budget and source
(Fig.~\ref{fig:rq2} shows 1,250),  ranging from \RQtwoRhoRange{}. The OHR-Bench documents of
held-out set~1 show the same pattern; on its \RQtwoHeldOutCUADDocuments{} CUAD
contracts the correlation weakens as the budget grows, and the one FinanceBench
document cannot be interpreted.

The pre-registered sign rule holds on held-out set~2 (Table~\ref{tab:rq2h2}): all
\RQtwoHTwoNegative{} of the \RQtwoHTwoValues{} values of $\rho$ are negative. Of
\RQtwoHTwoLargestBudgetValues{} intervals at each budget,
\RQtwoHTwoIncludeZeroSevenFiveZero{} includes zero at 750 tokens,
\RQtwoHTwoIncludeZeroOneTwoFiveZero{} at 1,250 and
\RQtwoHTwoLargestBudgetIncludeZero{} at 2,000. By corpus (pre-stated as
descriptive; primary gold), $\rho$ is positive at 2,000 tokens for HR
(\RQtwoHTwoHrDoclingTwoZeroZeroZero{} on Docling,
\RQtwoHTwoHrPdfplumberTwoZeroZeroZero{} on pdfplumber) and Industrial
(\RQtwoHTwoIndustrialDoclingTwoZeroZeroZero{} and
\RQtwoHTwoIndustrialPdfplumberTwoZeroZeroZero{}), and already at 1,250 for
Industrial on pdfplumber (\RQtwoHTwoIndustrialPdfplumberOneTwoFiveZero{}).

\subsection{RQ3: the measure changes the best size}\label{sec:rq3}

% Generated by paper_ieee/scripts/make_artifacts.py; do not edit.
\begin{figure}[!t]
\centering
\input{generated/figures/rq3.tex}
\caption{RQ3 on tuning: the same settings scored with a fixed number of chunks and with a fixed token budget. Both series use the questions of the fixed-k run (QASPER excluded; each parser's own gold), so Recall@1,250 here differs (up to \RQthreeFigureGap{}) from the tuning means on the questions kept for every setup; trend lines are least-squares fits for display only. Recall@4,000, a measure named in the diagnostic's design and outside the main budgets, also favours small chunks: best setups \RQthreeDoclingRecallFourZeroZeroZeroBest{} (\RQthreeDoclingRecallFourZeroZeroZeroSize{} tokens) on Docling and \RQthreePdfplumberRecallFourZeroZeroZeroBest{} (\RQthreePdfplumberRecallFourZeroZeroZeroSize{} tokens) on pdfplumber.}
\label{fig:rq3}
\end{figure}


On tuning data, when the number of chunks is fixed instead of the tokens, the
ranking reverses: on both parsers, Recall@10 rises with chunk size while
Recall@1,250 falls (Fig.~\ref{fig:rq3}). Recall@5 is highest for
\RQthreeDoclingRecallFiveBest{} on both parsers (measured size
\RQthreeDoclingRecallFiveSize{} and \RQthreePdfplumberRecallFiveSize{} tokens). MRR
and Hit@1 barely separate sizes; their winners are nominal, without intervals or
tests. In a control on tuning data, over four chunker families and two parsers,
size 256 minus size 1,024 is \RQthreeSweepBudget{} in Recall@1,250 but
\RQthreeSweepRecall{} in Recall@5; in Lift@5, where the recall of a random chunk
order is subtracted, it is \RQthreeSweepLift{}. All \RQthreeSweepSignificant{} of
the \RQthreeSweepTests{} tested differences are significant (largest Holm $p$ =
\RQthreeSweepMaxHolm{}). Under a fixed number of chunks, a large chunk puts more
text in front of the reader, so Recall@$k$ rewards size itself. 

\subsection{RQ4: retrieval quality versus answer quality}\label{sec:rq4}

% Generated by paper_ieee/scripts/make_artifacts.py; do not edit.
\begin{table}[!t]
\centering
\caption{RQ4 on held-out 1: token F1 of the best setup minus each setup.}
\label{tab:rq4}
\footnotesize
\begin{tabular}{lrrr}
\toprule
Setup & F1 difference & 95\% interval & Holm $p$ \\
\midrule
D/FT-1024 & +0.123 & [0.050, 0.210] & $<$0.01 \\
D/M1-384 & +0.065 & [0.010, 0.137] & 0.025 \\
D/LC-def & +0.048 & [0.002, 0.099] & 0.176 \\
D/LC-256 & +0.035 & [$-$0.021, 0.117] & 1.000 \\
D/Rec-512 & +0.030 & [$-$0.028, 0.113] & 1.000 \\
D/LC-512 & +0.010 & [$-$0.033, 0.051] & 1.000 \\
\bottomrule
\end{tabular}
\end{table}


On held-out set~1 (\AQHeldOutQuestions{} sampled questions), mean token F1 is
highest with \RQfourBest{} (\RQfourBestFOne{}; Table~\ref{tab:rq4}). Answers are
consistent with a large-chunk penalty, but this rests on one 1,024-token setup
that also differs in chunker. The best setup also beats
\RQfourDoclingMethodOneTreeTThreeEightFourSetup{}, although the two are tied on
recall; we report this as open and do not claim it. On tuning data
(\AQTuningQuestions{} questions), no difference is significant.  Across setups, Spearman's $\rho$ between F1
and Recall@1,250 is \RQfourSpearmanTuningFOne{} on tuning and
\RQfourSpearmanHeldOutFOne{} on held-out set~1. With only \RQfourSetups{} setups and
no test this is weak evidence, with a measure chosen after the answers existed.
On scanned PDFs (\AQScannedQuestions{} questions), the best setup is
\RQfourScannedBest{} (F1 \RQfourScannedBestFOne{}), with Docling's default OCR. Only one difference survives Holm correction, against
\RQfourScannedTopSetup{} (Holm $p$ = \RQfourScannedTopHolm{}). That pair differs in chunker as well as OCR, so this is
weak evidence on OCR.

\subsection{RQ5: tuning versus defaults}\label{sec:rq5}

\tuningfigure

\textbf{(a) Parser settings.} On digital PDFs (\RQfiveADigitalDocuments{} tuning
documents, \RQfiveADigitalQuestions{} shared questions), no parser variant differs
from its default by more than \RQfiveADigitalMaxAbsDiff{} in recall, and no
difference is significant (smallest Holm $p$ = \RQfiveADigitalMinHolm{};
Fig.~\ref{fig:rq5}a). On scanned PDFs (\RQfiveAScannedDocuments{} documents,
\RQfiveAScannedQuestions{} shared questions), we detect no difference between
Tesseract OCR and Docling's default OCR: \RQfiveAScannedDiff{} \RQfiveAScannedCI{}.
The parser choice that matters for scans is detecting them at all: pdfplumber
reads no text from scanned pages.

\textbf{(b) Chunk size.} In this post-hoc comparison on held-out set~1
(Fig.~\ref{fig:rq5}b), LangChain's default chunk size loses
\RQfiveBLangchainRecursiveDoclingDiff{} \RQfiveBLangchainRecursiveDoclingCI{} on
Docling and \RQfiveBLangchainRecursivePdfplumberDiff{}
\RQfiveBLangchainRecursivePdfplumberCI{} on pdfplumber, both with Holm $p$
\RQfiveBLangchainRecursiveDoclingHolm{}. For fixed-token, recursive and M1, whose
default is our own 512-token starting size, the tuned size differs from the
default by \RQfiveBRecursiveDoclingDiff{} to \RQfiveBMethodOneTreePdfplumberDiff{},
and no difference is significant.

\subsection{Extension, robustness and the tool}\label{sec:extension}

The NO-GO verdict is unchanged in descriptive checks. On held-out set~1, the gap is
\RobBMTwoFiveOnlyHeldOutOneGap{} \RobBMTwoFiveOnlyHeldOutOneGapCI{} with BM25 alone
and \RobDenseOnlyHeldOutOneGap{} \RobDenseOnlyHeldOutOneGapCI{} with dense retrieval
alone, and \RobWithoutIndianHeldOutOneGap{} \RobWithoutIndianHeldOutOneGapCI{}
without the Indian documents. On tuning data, adding the \ExtSetups{} setups of four
planned chunkers, including Docling's HybridChunker~\cite{hybridchunker}, leaves the
gap at \RobExtensionTuningGap{} \RobExtensionTuningGapCI{}. The HybridChunker result
includes our span mapping: \ExtHybridUnlocated{} of its \ExtHybridChunks{} chunks
could not be located, so it is not a test of HybridChunker as used directly. With
nomic-embed-text-v1.5~\cite{nomic} as the dense embedder, also on tuning data, the
gap is \RobNomicHybridTuningGap{} \RobNomicHybridTuningGapCI{} with hybrid
retrieval. The robustness table and the RQ3 erratum are in the repository (tag paper-v1.4).

The tool does not choose a setup per document. Its default mode applies
\RQoneTuningWinner{} to every PDF; its fast mode keeps the chunker but parses with
pdfplumber, unless a page has no text layer; then the whole document goes to
Docling. Fast mode on digital PDFs is therefore not the tuned setup.


\section{Discussion}\label{sec:discussion}

\textbf{What the NO-GO means.} Neither held-out set gives the evidence the decision rule asks
for, that a choice per document beats the best single setup.
This is not proof that a better chooser can never help: the best single setup is
selected on the same questions it is scored on, and the oracle chooses from only
half of a document's few questions.  Within the cases tested (English PDFs, two parsers, 1,250
tokens), this supports a simple default: one small-chunk setup rather than a
per-document selector.

\textbf{Fixed $k$ versus fixed budget.} RQ2 and RQ3 (the latter on tuning data)
together suggest that the best chunk size is not a property of the chunker alone.
Under a fixed number of chunks, a larger chunk shows the reader more text, so
Recall@$k$ rewards size itself; under a fixed budget of 750 or 1,250 tokens that
advantage disappears and small chunks tend to win. Comparisons at a fixed number of chunks therefore mix chunk quality with the
amount of text read; we suggest that they report the token budget, or both
measures.

\textbf{Defaults.} Tuning paid off only against LangChain's default chunk size
(post-hoc, Section~\ref{sec:rq5}). For settings like those tested, a practical check is to
compare a library's default chunk size with the token budget in use and, when a
parser without OCR is used, to detect scanned pages.


\section{Limitations}\label{sec:limitations}

All scored documents and questions are in English. The main results use one
embedder, bge-m3; the second was run on tuning data only.  RQ1 was analysed at 1,250 tokens only; the planned analyses at 750 and 2,000
were not run. Many
questions were generated by language models, and some Indian ones are invalid
(Section~\ref{sec:docs}); the Indian held-out questions were not checked.
Held-out set~1 was used more than once, and the gold source and rule differ
between the two sets; on held-out set~2 the drop rule replaced the earlier-fixed
certain-only rule after gold counts, though before any score. Held-out set~2
excludes questions on charts and images and documents over \HTwoMaxPages{}
pages; its box gold is read from the PDF text layer, which pdfplumber also
extracts from, so pages without a text layer give no gold. RQ4 uses one
answering model, and \AQUnrecordedAnswers{} of \AQAnswers{} answers have no run
record. 


\section{Conclusion}\label{sec:conclusion}

On two sealed held-out sets, at 1,250 tokens and under a decision rule written in
advance, a cross-fitted per-document oracle did not beat the best single setup,
and we find no evidence of a gain from per-document choice. This is not a bound
on every chooser. Small chunks tend to retrieve best at a budget of 750 or 1,250
tokens, which is not established at 2,000; on tuning data the ranking reverses
when the number of chunks is fixed instead; and in a post-hoc comparison, tuning
helps against LangChain's default chunk size.

\section*{Data and Code Availability}
Code, specifications, run records, checksum files, and score files and result
tables that hold identifiers, scores and statistics only are at
\url{https://github.com/amanchaurasia512/Self_optimized_Agent_Ai_chunking}, tag
\texttt{paper-v1.4}. The project's own code is under the MIT licence;
\texttt{NOTICE.md} lists the terms of material derived from the datasets. CUAD,
QASPER, the questions of OHR-Bench and the annotations of ViDoRe~V3 are under CC~BY~4.0. UniDoc-Bench and the annotations of FinanceBench are under
CC~BY-NC~4.0. MP-DocVQA is under MIT at its source. The repository of the Indian
benchmark is under Apache-2.0. Document licences vary or are unknown. The PDFs, extracted text,
chunks, embeddings and model answers are not redistributed.

\section*{Acknowledgment}
Following IEEE policy, I disclose my use of AI in this work. I used Claude
(Anthropic), through the Claude app and Claude Code, between October and
November 2026. Claude drafted the text of the Abstract, Sections~I--VII and the data and code availability statement, working from my study
design, my dated specifications and the locked results, and condensed the full
paper into this six-page version, which I checked. For Related Work, it
drafted from the abstracts of the cited papers, and each statement was
checked against them. Claude Code wrote most of the analysis code and the
script that produces every table, figure and reported number. No figure or
image was made by an AI model. I designed the study, made every research
decision and ran the experiments. I reviewed and edited every section,
checked the outputs of all the code, and I take full responsibility for the
content of this paper.
\bibliographystyle{IEEEtran}
\makeatletter\let\input@path\undefined\makeatother % main.bbl is this paper's own
\bibliography{bstcontrol,../paper_ieee/references}

\end{document}
